\section{Introduction}\label{sec:intro}
In object-oriented languages, \textit{method overloading} and its
unique resolution are important features.  Method overloading allows
multiple method declarations of the same name but different parameter
types.  Thus, it enables to define different behaviors of $+$ or
\texttt{equal} for different parameter types, for example.  For a given
overloaded method invocation of \texttt{equal}, a resolution mechanism
should select a single declaration to call among the overloaded
declarations of \texttt{equal} using the argument types of the method invocation.

When resolving an overloaded method invocation, \textit{dynamic dispatch}
considers the run-time types of the arguments of the method invocation.
When a dispatch mechanism considers the run-time type of only a single
argument, usually the receiver of a method call, as in languages like
\CC and Java, it is \textit{single dispatch}.  By contrast,
\textit{multiple dispatch} considers the run-time types of more than one
argument as in languages like Common Lisp, Fortress, and Julia.
Especially, \textit{symmetric multiple dispatch} considers multiple
arguments equally, which supports mathematical operators like binary
methods intuitively and uniformly~\cite{clifton00}.

While method overloading and its resolution by dynamic dispatch
give a consistency in naming methods that logically perform very
similar tasks, they impose two restrictions to guarantee a unique
and type-sound resolution of each overloaded method
invocation~\cite{castagna92,castagna95}.  First, the dispatch mechanism should be
able to find the \textit{most specific applicable} method declaration for
any overloaded method invocation.  A method declaration is applicable to
given arguments if its parameter types are supertypes of the types of
the arguments.  One method declaration is more specific than another
if the \emph{domain type} of the former is a subtype of that of the
latter.  Hereafter, we consider that every method declaration has a
single parameter possibly a tuple of two or more components, and call the
parameter type of a method declaration its domain type.  Second, for a
given overloaded method invocation, its static type should be a supertype
of the return type of a dynamically chosen method declaration by the
dispatch mechanism.  Otherwise, a method invocation result may be used
in a type-unsound manner.

Moreover, expressive language features often add various possibilities to
make overloaded method invocations ambiguous.  As we discuss below,
such features make it difficult to design a language that guarantees no
ambiguous nor undefined calls at run time.  

One such a feature is \textit{multiple inheritance}, which allows one
class to extend multiple classes.  While multiple inheritance supports
modular and reusable code, it easily introduces ambiguous method
calls.  Because one parent class may not be more specific than another
parent class, dynamically selecting the most specific applicable method
declaration may be unclear.  For example, consider that each of two
classes $A$ and $B$ has a method declaration $m$ with the same
parameter type \texttt{Number}, class $C$ extends both $A$ and $B$, and
$c$ is an instance of $C$.  Then, evaluation of $c.m(0)$ results in an
ambiguous call because neither $A$ nor $B$ is more specific than each
other.

Another feature is \textit{parametric polymorphism}, or \textit{generics}.
By introducing type parameters in class and method declarations,
parametric polymorphism enables programmers to manage types in a
fine-grained way.
% For example, the \texttt{List} class may be refined as
% \texttt{List$\ob{\texttt{String}}$} and
% \texttt{List$\ob{\texttt{Integer}}$}, which allows different
% operations on lists of strings and lists of integers.
For example, the \texttt{List} class may be instantiated as either
\texttt{List$\ob{\texttt{String}}$} or
\texttt{List$\ob{\texttt{Integer}}$}.
However,
% parametric polymorphism also easily introduces ambiguity in overloading
% resolution.  Because polymorphic method declarations have domain types that
% are more generic than monomorphic method declarations,
selecting the most
specific applicable declaration among polymorphic method declarations
at run time is not trivial.
% Moreover, since languages with parametric polymorphism often support type
% inference for the type arguments of polymorphic method invocations, the
% dynamic dispatch should not only find the most specific applicable
% declaration, but also infer the type arguments of method invocations.
Moreover, many languages with parametric polymorphism support automatic
inference for the type arguments of polymorphic method invocations;
at run time, dynamic dispatch must not only find the most specific applicable
declaration, but also infer appropriate type arguments for that declaration.

Finally, \textit{variance} defines additional subtype relations for
polymorphic types~\cite{variance}.
% While polymorphic types without
% variance have only a restricted pointwise subtype relation, where
Without variance,
instances of the same polymorphic type are in a subtype relation only
when their type arguments are the same type.
Variance supports subtype
relationships between different instances of the same polymorphic type
by allowing the programmer to label each type parameter of a class
\textit{covariant}, \textit{contravariant}, or \textit{invariant}.
% ; more
% specifically, programmers can define class declarations with
% variance annotations of type parameters: \textit{covariant},
% \textit{contravariant}, and \textit{invariant}.
If a parametric type $T$
has one covariant type parameter, then it preserves the subtype relation
of its type argument; hence, when a type $P$ is a subtype of another type
$Q$, which we write $P\extends Q$, it implies $T\ob{P}\extends T\ob{Q}$.
If the type parameter of $T$ is contravariant, then $T$ reverses the
subtype relationship of its type arguments; hence, $P\extends Q$ implies
$T\ob{Q}\extends T\ob{P}$.  If the type parameter of $T$ is invariant,
then
% $T$ ignores any subtype relation between its type arguments; hence,
$T\ob{P}\extends T\ob{Q}$ only if $P\equiv Q$, where $P\equiv Q$ denotes
that $P$ and $Q$ are equivalent, that is, $P$ is a subtype of $Q$ and
vice versa.
% With such expressive subtype relations between polymorphic
% types, programmer can define more expressive polymorphic methods than
% without variance.
Variance annotations allow the programmer to define more expressive polymorphic methods.
For example, consider the following excerpt from the
\texttt{List} class source code in Scala~\cite{scalalist}:\\[-.7em]

{\tt\setlength{\baselineskip}{.5\baselineskip}\noindent
\hspace*{5em}sealed abstract class List[+A] extends Seq[A] \{\\
\hspace*{7em}$\cdots$\\[-.4em]
\hspace*{7em}def :::[B >: A](prefix: List[B]): List[B] = $\cdots$\\[-.4em]
\hspace*{7em}$\cdots$\\[-.4em]
\hspace*{5em}\}
}
\\[-.7em]

\noindent
The \texttt{List} class takes a single covariant (denoted by \texttt{+}) type parameter
\texttt{A}, and the
\texttt{:::} method concatenates two lists, its receiver and argument.
If \texttt{A} were invariant, the element type \texttt{B} of the argument
list should be a supertype of the element type \texttt{A} of the receiver
list.  However, since \texttt{A} is covariant, the method can take a list
of any element type \texttt{C} that is less than or equal to \texttt{B},
and return a list whose element type is
\texttt{B}, which is a common supertype of the element types of its
receiver and argument: \texttt{A}$\extends$\texttt{B} and
\texttt{C}$\extends$\texttt{B}.
% As expected, due to the increased
% subtype relations between polymorphic types, variance
% also adds ambiguities in overloading resolution.
As one might expect, variance also
increases the difficulty of overloading resolution.

Therefore, supporting symmetric multiple dispatch in the presence of
various expressive language features has been a long-standing effort.
Previous efforts have formalized a variety of core calculi that capture
some of the aforementioned features in a piecemeal rather than holistic
fashion.  While they are widely different in terms of the expressivity,
completeness, decidability, and correctness of the formalization of the
target languages and their type systems, they all define a set of
\emph{overloading rules} that overloaded method declarations should
satisfy to make dynamic dispatch select the most specific applicable
method declaration with a type-safe return type.

For a simple object-oriented language with multiple inheritance,
\citet{ff, fortress} and \citet{ffcoq} presented three overloading rules:
Exclusion Rule, Subtype Rule, and Meet Rule.
% Each rule gives constraints
% for a pair of two method declarations, and every pair of method
% declarations with the same name should satisfy \emph{one} of the
% overloading rules.
Every pair of method
declarations with the same name must satisfy \emph{one} of these
overloading rules.
The authors proved that a set of overloaded method
declarations satisfying the overloading rules does not produce any
ambiguous method invocation at run time.

To understand the interplay between symmetric multiple dispatch and
parametric polymorphism, \citet{oopsla11} focused on a type system
without any specific target language and informally presented three other
overloading rules: a No Duplicates Rule, a revised Meet Rule, and a Return Type Rule,
% As before, each rule puts constraints for a pair of two method
% declarations, but, this time, every pair of method declarations with the
% same name should satisfy \emph{all} three overloading rules.
where every pair of method declarations with the
same name must satisfy \emph{all three} overloading rules.
The authors
showed that interpreting a polymorphic method declaration as a set of
infinitely many monomorphic declarations is not appropriate for
overloading resolution.  Instead, they considered a polymorphic method
declaration as a single declaration whose domain type is existentially
quantified over its type parameters as proposed by \citet{bm97}.

%Based on the previous work, we finally
% present a formal specification of
% an object-oriented language with symmetric multiple dispatch, multiple
% inheritance, and parametric polymorphism with variance.
% We define a core calculus named \emph{FGFV}, Featherweight Generic Fortress with Variance,
% which contains only the core language features.
% For example, FGFV omits functional
% methods~\cite{ff} and declarations of closed types~\cite{oopsla11}
% (but these could be added to FGFV in an obvious way).
% We formally present overloading
% rules for FGFV and define its dynamic semantics with a new dynamic
% dispatch algorithm.  We also prove the type soundness of FGFV.

In this paper, we
present a formal core calculus \emph{FGFV} (Featherweight Generic Fortress with Variance)
for an object-oriented language with symmetric multiple dispatch, multiple
inheritance, and parametric polymorphism with variance.
For conciseness, FGFV omits functional methods~\cite{ff} and declarations of closed types~\cite{oopsla11},
but these could be added back to FGFV in an obvious way.
Unlike the prior work, we formally present overloading rules for FGFV,
define its dynamic semantics with a novel dynamic
dispatch algorithm, and prove type soundness.
As we discuss in Section~\ref{sec:discuss}, defining a provable system
for an object-oriented language with asymmetric features requires careful
tradeoffs between static expressivity and dynamic efficiency.

In the rest of the paper, we first describe the previous efforts in more
detail with concrete code examples (Section~\ref{sec:motiv}).  After
presenting the FGFV calculus that represents an object-oriented
language with symmetric multiple dispatch, multiple inheritance,
parametric polymorphism, and variance (Section~\ref{sec:lang}), we
formally define the overloading rules for FGFV (Section~\ref{sec:rules}).
Then, we define both static and dynamic overloading resolution with
a novel dynamic dispatch algorithm, and prove the type soundness of FGFV
(Section~\ref{sec:resolve}).  Finally, we discuss related work
(Section~\ref{sec:related}) and conclude (Section~\ref{sec:conc}).
